\documentclass[12pt]{article}
\usepackage[brazilian]{babel}
\usepackage[
  perfil=academico,
  impressao=frente-e-verso
]{abntex3}

\ABNTEXsetup{
  titulo={Elementos pré-textuais no abnTeX3 Community},
  subtitulo={exemplo integrado do Marco 5},
  autoria={Maria da Silva},
  instituicao={Universidade Exemplo},
  natureza={Dissertação},
  objetivo={apresentada para obtenção do título de Mestre},
  area={Documentação},
  orientacao={Prof. Dr. João Souza},
  coorientacao={Profa. Dra. Ana Santos},
  local={Recife},
  data={2026}
}

\ABNTEXabreviatura{Fil.}{Filosofia}
\ABNTEXsigla{ABNT}{Associação Brasileira de Normas Técnicas}
\ABNTEXsigla{IBGE}{Instituto Brasileiro de Geografia e Estatística}
\ABNTEXsimbolo{$n$}{Quantidade de elementos}
\ABNTEXsimbolo{$\sum$}{Somatório}

\begin{document}

\ABNTEXcapa
\ABNTEXlombada[Dissertação 2026][Universidade Exemplo]

\ABNTEXpretextualpages
\ABNTEXfolhaderosto

\ABNTEXfichacatalografica{%
  Dados internacionais de catalogação na publicação devem ser preparados
  conforme o código de catalogação vigente e inseridos neste ponto de
  extensão pelo profissional responsável.
}

\ABNTEXerrata
  {SILVA, Maria da. \textbf{Elementos pré-textuais no abnTeX3 Community}.
   2026. Dissertação — Universidade Exemplo, Recife, 2026.}
  {\begin{center}
     \begin{tabular}{llll}
       Folha & Linha & Onde se lê & Leia-se\\
       10 & 4 & texo & texto
     \end{tabular}
   \end{center}}

\ABNTEXfolhadeaprovacao
  {28 de julho de 2026}
  {%
    \ABNTEXmembrobanca
      {Prof. Dr. João Souza}{Orientador}{Universidade Exemplo}
    \ABNTEXmembrobanca
      {Profa. Dra. Ana Santos}{Examinadora}{Universidade Exemplo}
  }

\ABNTEXdedicatoria{À comunidade brasileira de usuários de \LaTeX.}

\ABNTEXagradecimentos{%
  A todas as pessoas que contribuíram para a construção e a revisão deste
  exemplo.
}

\ABNTEXepigrafe[Autoria demonstrativa]{%
  A documentação também faz parte do programa.
}

\ABNTEXresumo
  {Este exemplo demonstra a composição dos elementos pré-textuais, o
   registro das listas e a transição para a parte textual por meio da API
   pública do abnTeX3 Community.}
  {normalização, documentação, LaTeX}

\ABNTEXresumoemlinguaestrangeira
  {This example demonstrates front matter composition, list registration,
   and the transition to the textual part through the public API.}
  {standardization, documentation, LaTeX}

\ABNTEXlistadeilustracoes
\ABNTEXlistadetabelas
\ABNTEXlistadeabreviaturas
\ABNTEXlistadesiglas
\ABNTEXlistadesimbolos
\ABNTEXsumario

\ABNTEXtextualpages

\section{Introdução}

O conteúdo textual começa com numeração visível e preserva a contagem das
páginas pré-textuais.

\begin{figure}[ht]
  \centering
  \rule{0.6\linewidth}{2cm}
  \caption{Demonstração de uma ilustração}
  \ABNTEXfonte{elaborada pela própria autora.}
\end{figure}

\begin{table}[ht]
  \centering
  \caption{Demonstração de uma tabela}
  \begin{tabular}{ll}
    Elemento & Estado\\
    Capa & implementado\\
    Resumo & implementado
  \end{tabular}
  \ABNTEXfonte{elaborada pela própria autora.}
\end{table}

\end{document}
